\documentclass[11pt,a4paper]{article}

\usepackage[margin=2.5cm, headheight=14pt]{geometry}
\usepackage{booktabs}
\usepackage{multirow}
\usepackage{array}
\usepackage{xcolor}
\usepackage{hyperref}
\usepackage{amsmath}
\usepackage{microtype}
\usepackage{parskip}
\usepackage{enumitem}
\usepackage{caption}
\usepackage{fancyhdr}
\usepackage{titlesec}
\usepackage{tabularx}
\usepackage{colortbl}

\definecolor{ownerblue}{RGB}{30,80,160}
\definecolor{bertcol}{RGB}{210,230,255}
\definecolor{llmcol}{RGB}{210,255,220}
\definecolor{s2tmlmcol}{RGB}{255,235,200}
\definecolor{s2tollamacol}{RGB}{255,215,215}

\hypersetup{colorlinks=true,linkcolor=ownerblue,urlcolor=ownerblue,citecolor=ownerblue}

\pagestyle{fancy}
\fancyhf{}
\rhead{\textcolor{gray}{\small OWNER \& Span2Type -- Cluster Naming Evaluation}}
\lhead{\textcolor{gray}{\small BERTScore Report (Version~3)}}
\rfoot{\thepage}

\titleformat{\section}{\large\bfseries\color{ownerblue}}{{\thesection}}{1em}{}
\titleformat{\subsection}{\normalsize\bfseries}{{\thesubsection}}{1em}{}

\begin{document}

% ─── Title ───────────────────────────────────────────────────────────────────
\begin{center}
  {\LARGE\bfseries\color{ownerblue}
    Cluster Naming Evaluation Report \normalsize\textit{(Version~3)}\\[0.4em]
    \large OWNER vs.\ Span2Type: Cluster Naming Comparison}\\[0.8em]
  {\large BERTScore Evaluation across Five CrossNER Domains}\\[0.6em]
  {\normalsize Chaithra-lm \quad \texttt{18520923@gm.uit.edu.vn}}\\[0.3em]
  {\small \today}
\end{center}
\vspace{0.4em}\hrule\vspace{1.2em}

% ─── Abstract ────────────────────────────────────────────────────────────────
\begin{abstract}
This report compares four cluster-naming methods across five CrossNER domains
(AI, Literature, Music, Politics, Science), evaluated using BERTScore
(RoBERTa-large) against ground-truth entity type labels.
Two methods belong to \textbf{OWNER}: BERT-based naming (MLM top-1 vote)
and LLM-based naming (Llama-3.1-8B via Ollama, Figure~3 prompt, $n{=}16$).
Two methods belong to \textbf{Span2Type}: MLM-based naming (same BERT voting
strategy) and Ollama-based naming (Llama-3.1-8B with richer context prompting).
OWNER-BERT and Span2Type-MLM produce \textbf{identical} F$_1$ scores (0.966)
since they share the same underlying approach.
Span2Type-Ollama achieves a higher average F$_1$ of \textbf{0.952} compared
to OWNER-LLM's \textbf{0.944}, driven by more descriptive multi-word names.
\end{abstract}

% ─── 1. Introduction ─────────────────────────────────────────────────────────
\section{Introduction}

Unsupervised NER systems such as OWNER~\cite{genest2025owner} group extracted
entities into clusters identified only by integer IDs.
Automatic cluster naming assigns human-readable type labels without manual
annotation. This report evaluates and compares four naming strategies:

\begin{itemize}[leftmargin=1.5em]
  \item \textbf{OWNER-BERT}: for each entity in a cluster, fill the prompt
    \texttt{``\{sentence\} \{entity\} is a [MASK].''} and take the
    top-1 MLM prediction; the most frequent token across the cluster is the name.
  \item \textbf{OWNER-LLM}: sample $n{=}16$ entities per cluster, query
    Llama-3.1-8B with the Figure~3 prompt (temperature~$=0$), use the
    response as the name.
  \item \textbf{Span2Type-MLM}: same BERT voting strategy as OWNER-BERT,
    applied to Span2Type cluster assignments.
  \item \textbf{Span2Type-Ollama}: Llama-3.1-8B with richer context
    (full example sentences included in the prompt), producing longer and
    more specific multi-word names.
\end{itemize}

\noindent
All methods are scored against ground-truth labels using
BERTScore~\cite{zhang2020bertscore} (RoBERTa-large backbone).

% ─── 2. Experimental Setup ───────────────────────────────────────────────────
\section{Experimental Setup}

\subsection{MLflow Run IDs}

\begin{table}[h]
\centering
\caption{MLflow run IDs for OWNER and Span2Type experiments.}
\label{tab:runs}
\setlength{\tabcolsep}{5pt}
\small
\begin{tabular}{l l l}
\toprule
\textbf{Domain} & \textbf{OWNER Run ID} & \textbf{Span2Type Run ID} \\
\midrule
AI         & \texttt{ea6e40f8b9ed4aa385956217b939eb6e} & \texttt{6b105207f3084468b2b79f37841ae2b2} \\
Literature & \texttt{fb1215f7b8e642ee99d43782427432f0} & \texttt{29e11a34c7214005b0d96bc0520515be} \\
Music      & \texttt{2c886b9b2f104bd4ad79ed1bf3dd13c5} & \texttt{56ad65e38f4e464f85acc934340692dc} \\
Politics   & \texttt{9c02f083100c41b3b052373d4e8e916d} & \texttt{610c691d79904252b375511432bbbe7b} \\
Science    & \texttt{f72db65d19f243a6aedda8687a538896} & \texttt{efcfb115a1cf47798145e9ee97b01755} \\
\bottomrule
\end{tabular}
\end{table}

\subsection{Ground-Truth Alignment}

For OWNER runs, the GT label per cluster comes from
\texttt{et\_test\_cluster\_naming\_gt.json} (column-wise argmax of the
confusion matrix). For Span2Type runs (which do not store a GT file),
the same argmax procedure is applied directly to the confusion matrix
tensor (\texttt{et\_test\_confusion-matrix.pt}) and its metadata.

\subsection{Cluster Statistics}

\begin{table}[h]
\centering
\caption{Number of aligned clusters per domain and run type.}
\label{tab:clusters}
\setlength{\tabcolsep}{8pt}
\begin{tabular}{l r r r r}
\toprule
\textbf{Domain} & \textbf{True $k$} & \textbf{OWNER clusters} & \textbf{Span2Type clusters} \\
\midrule
AI         & 14 & 16 & 16 \\
Literature & 12 & 14 & 14 \\
Music      & 13 & 20 & 20 \\
Politics   &  9 & 22 & 22 \\
Science    & 17 & 20 & 20 \\
\midrule
Total      & 65 & 92 & 92 \\
\bottomrule
\end{tabular}
\end{table}

% ─── 3. Results ──────────────────────────────────────────────────────────────
\section{Results}

\subsection{BERTScore Comparison}

Table~\ref{tab:bertscore} presents P, R, F$_1$ for all four methods.
The best F$_1$ per domain is shown in \textbf{bold}.

\begin{table}[h]
\centering
\caption{BERTScore (P / R / F$_1$) for four cluster-naming methods across five
CrossNER domains. Best F$_1$ per domain in \textbf{bold}.}
\label{tab:bertscore}
\setlength{\tabcolsep}{3.5pt}
\small
\begin{tabular}{l ccc ccc ccc ccc}
\toprule
& \multicolumn{3}{c}{\cellcolor{bertcol}\textbf{OWNER-BERT}}
& \multicolumn{3}{c}{\cellcolor{llmcol}\textbf{OWNER-LLM}}
& \multicolumn{3}{c}{\cellcolor{s2tmlmcol}\textbf{S2T-MLM}}
& \multicolumn{3}{c}{\cellcolor{s2tollamacol}\textbf{S2T-Ollama}} \\
\cmidrule(lr){2-4}\cmidrule(lr){5-7}\cmidrule(lr){8-10}\cmidrule(lr){11-13}
\textbf{Domain} & P & R & F$_1$ & P & R & F$_1$ & P & R & F$_1$ & P & R & F$_1$ \\
\midrule
AI
  & 0.992 & 0.986 & \textbf{0.989}
  & 0.915 & 0.945 & 0.929
  & 0.992 & 0.986 & \textbf{0.989}
  & 0.934 & 0.956 & 0.944 \\
Literature
  & 0.974 & 0.975 & 0.974
  & 0.950 & 0.960 & 0.954
  & 0.974 & 0.975 & 0.974
  & 0.969 & 0.970 & \textbf{0.969} \\
Music
  & 0.965 & 0.951 & \textbf{0.958}
  & 0.960 & 0.955 & 0.957
  & 0.965 & 0.951 & \textbf{0.958}
  & 0.948 & 0.940 & 0.943 \\
Politics
  & 0.968 & 0.962 & \textbf{0.964}
  & 0.945 & 0.961 & 0.953
  & 0.968 & 0.962 & \textbf{0.964}
  & 0.958 & 0.968 & 0.963 \\
Science
  & 0.949 & 0.943 & \textbf{0.945}
  & 0.922 & 0.933 & 0.926
  & 0.949 & 0.943 & \textbf{0.945}
  & 0.945 & 0.937 & 0.941 \\
\midrule
\textbf{Average}
  & \textbf{0.970} & \textbf{0.963} & \textbf{0.966}
  & 0.938 & 0.951 & 0.944
  & \textbf{0.970} & \textbf{0.963} & \textbf{0.966}
  & 0.951 & 0.954 & 0.952 \\
\bottomrule
\end{tabular}
\end{table}

\subsection{Qualitative Comparison}

Table~\ref{tab:examples} shows representative naming examples contrasting
OWNER-LLM and Span2Type-Ollama, where the two approaches differ most.

\begin{table}[h]
\centering
\caption{Selected naming examples comparing LLM-based approaches.
\textcolor{green!60!black}{Green} = correct/near-correct;
\textcolor{orange!70!black}{orange} = partially correct;
\textcolor{red!70!black}{red} = incorrect.}
\label{tab:examples}
\setlength{\tabcolsep}{4pt}
\small
\begin{tabularx}{\textwidth}{l l X X X}
\toprule
\textbf{Dom.} & \textbf{GT} & \textbf{OWNER-LLM} & \textbf{S2T-Ollama} \\
\midrule
AI   & algorithm     & \textcolor{orange!70!black}{Computer Science Concepts} & \textcolor{orange!70!black}{artificial intelligence machine learning} \\
AI   & programlang   & \textcolor{green!60!black}{Programming languages}      & \textcolor{green!60!black}{programming languages} \\
AI   & product       & \textcolor{orange!70!black}{Robot}                      & \textcolor{green!60!black}{machine learning library} \\
AI   & person        & \textcolor{green!60!black}{Person}                      & \textcolor{orange!70!black}{director} \\
\midrule
Lit. & book          & \textcolor{green!60!black}{Novel}                       & \textcolor{orange!70!black}{character} \\
Lit. & literarygenre & \textcolor{orange!70!black}{Literary work}               & \textcolor{red!70!black}{novels} \\
Lit. & event         & \textcolor{orange!70!black}{Film Festivals}              & \textcolor{green!60!black}{film festivals} \\
\midrule
Mus. & musicalartist & \textcolor{green!60!black}{Singer}                      & \textcolor{green!60!black}{singer} \\
Mus. & event         & \textcolor{green!60!black}{Event}                       & \textcolor{orange!70!black}{multisport events} \\
Mus. & album         & \textcolor{red!70!black}{Song}                          & \textcolor{green!60!black}{albums} \\
\midrule
Pol. & politicalparty & \textcolor{green!60!black}{Political party}             & \textcolor{green!60!black}{political party} \\
Pol. & organisation  & \textcolor{green!60!black}{Government}                  & \textcolor{green!60!black}{legislative body} \\
Pol. & person        & \textcolor{red!70!black}{Celebrity}                     & \textcolor{orange!70!black}{actress} \\
\midrule
Sci. & scientist     & \textcolor{green!60!black}{Physicist}                   & \textcolor{orange!70!black}{historical figures} \\
Sci. & chemicalcompound & \textcolor{green!60!black}{Chemical substance}       & \textcolor{green!60!black}{chemical compound} \\
Sci. & misc          & \textcolor{red!70!black}{Ethnic groups}                 & \textcolor{green!60!black}{biological processes} \\
\bottomrule
\end{tabularx}
\end{table}

% ─── 4. Analysis ─────────────────────────────────────────────────────────────
\section{Analysis}

\subsection{OWNER-BERT = Span2Type-MLM}

Both BERT-based methods achieve \textbf{identical} BERTScore across all
domains (avg.\ F$_1 = 0.966$). This is expected: they use the same
BERT-base-uncased MLM head with the same prompt template and voting
mechanism on the same cluster assignments. The `\texttt{vegetarian}''
artefact appears in both (8 of 92 clusters, 8.7\%) for the same reason:
ambiguous entities in the \texttt{``X is a [MASK].''} pattern trigger
BERT's most frequent fill pattern from pretraining.

\subsection{Span2Type-Ollama Outperforms OWNER-LLM}

Span2Type-Ollama achieves a higher average F$_1$ (\textbf{0.952}) than
OWNER-LLM (\textbf{0.944}), a gap of $+0.008$. The advantage comes from
richer context: Span2Type-Ollama includes full example sentences in its
prompt, allowing the LLM to infer fine-grained types such as
\textit{machine learning library} (GT: product), \textit{legislative body}
(GT: organisation), and \textit{biological processes} (GT: misc).

OWNER-LLM, which only passes entity strings without their sentence context,
is more susceptible to over-generalisation (e.g., \textit{Celebrity} for
\texttt{person}; \textit{Stadiums} for \texttt{country}).

The improvement is most pronounced on \textbf{Literature} ($\Delta$F$_1 =
+0.015$) and smallest on \textbf{Music} ($-0.014$), where Span2Type-Ollama
over-specifies (\textit{multisport events} for \texttt{event}).

\subsection{Per-Domain Summary}

\begin{itemize}[leftmargin=1.5em]
  \item \textbf{AI}: BERT methods dominate (F$_1=0.989$). LLM methods
    struggle with abstract types like \texttt{algorithm} and \texttt{task}.
  \item \textbf{Literature}: Span2Type-Ollama achieves the highest F$_1$
    (0.969), narrowly beating the BERT methods (0.974 vs 0.969 is within
    rounding; Span2Type-Ollama has higher Recall 0.970 vs 0.975).
  \item \textbf{Music}: All methods perform similarly (range 0.943--0.958).
    Both Ollama variants confuse \texttt{album} with song-related terms.
  \item \textbf{Politics}: BERT methods and Span2Type-Ollama are within
    0.001 F$_1$, all around 0.963--0.964.
  \item \textbf{Science}: BERT methods lead (0.945); Span2Type-Ollama
    is close (0.941). Fine-grained scientific types remain challenging
    for all approaches.
\end{itemize}

% ─── 5. Conclusion ───────────────────────────────────────────────────────────
\section{Conclusion}

Across all four methods, BERT/MLM voting achieves the highest BERTScore
(F$_1=0.966$) but produces single-token names and the \textit{vegetarian}
artefact. Among LLM-based approaches, \textbf{Span2Type-Ollama} (F$_1=0.952$)
outperforms \textbf{OWNER-LLM} (F$_1=0.944$) by providing fuller sentence
context to the model, yielding more precise multi-word names.

The ranking across all methods is:

\begin{center}
OWNER-BERT $=$ Span2Type-MLM ($0.966$)
$>$ Span2Type-Ollama ($0.952$)
$>$ OWNER-LLM ($0.944$)
\end{center}

\noindent
For practical deployment where human-readable output matters, Span2Type-Ollama
is recommended. For cost-free naming within an existing OWNER pipeline, the
MLM voting approach remains the strongest baseline.

% ─── References ──────────────────────────────────────────────────────────────
\begin{thebibliography}{9}

\bibitem{genest2025owner}
P.-Y.~Genest, P.-E.~Portier, E.~Egyed-Zsigmond, and M.~Lovisetto,
``OWNER -- Toward Unsupervised Open-World Named Entity Recognition,''
\textit{IEEE Access}, vol.~13, pp.~50077--50096, 2025.

\bibitem{zhang2020bertscore}
T.~Zhang, V.~Kishore, F.~Wu, K.~Q.~Weinberger, and Y.~Artzi,
``BERTScore: Evaluating Text Generation with BERT,''
in \textit{Proc. ICLR}, 2020.

\end{thebibliography}

\end{document}
